\begin{abstract}
Modular robots built by a mirror-symmetric genome are often reported to walk further than robots built by an unconstrained genome. We ask whether mirroring is the reason. On fixed random bodies with equal hinge counts, symmetric bodies walk 18--25\% further than unconstrained ones once their controllers are trained, a modest effect that is clearest with an open-loop controller. The lead is absent at the start of training and builds during learning. Two kinds of unmirrored bodies match symmetric bodies: unconstrained bodies selected for the same resting shape, and bodies that duplicate modules without mirroring them. In evolution, bodies that duplicate modules also end ahead of unconstrained ones, while the symmetric lead is similar in size but not significant over 20 seeds. Symmetric bodies do move differently, locking their left and right limbs in phase, and locked gaits score higher. But perturbing the left--right timing of a trained gait costs no more than perturbing any other timing, so trained gaits do not depend on the lock. Mirroring is therefore not required for the symmetry advantage, and comparisons of body encodings should control for the body shape and duplication an encoding produces.
\end{abstract}
